\section{Conclusion}
\label{sec:conclusion}

A multi-provider lunar PNT architecture needs a number that no interoperability
standard yet publishes: the consistency tolerance, the maximum cross-provider
frame-and-dynamics disagreement a user can absorb within a positioning budget. The
LunaNet Interoperability Specification defers it to a still-unreleased applicable
document---AD-5, the Lunar Reference System and LunaNet Reference Time System Standard
({LNIS-TBD-AD0005})---which does not yet specify one. We have supplied it in three
parts. A consistency-tolerance theorem shows the interop tax is a free-network gauge
phenomenon---a datum common to all providers is invisible to a single provider and
cancels in provider-versus-provider fusion, while the \emph{differential} datum that
two providers do not share is the tax---that splits, as a measured affine
decomposition of the recovered rotation, into a reducible part a common frame removes
and an irreducible part only a common ephemeris removes, and that inverts into a
closed-form tolerance $\tau(B)=B/R_{\mathrm{Moon}}$, binding on orientation, for a
user budget $B$. The only real ``multiple providers'' that exist today---the three
independent authoritative lunar ephemerides DE440, INPOP21a, and EPM2021---anchor the
budget: their measured 2024--2025 geocentric-Moon disagreement is an orientation
effect (a constant rotation explains $71$--$94\%$ by amplitude, $91$--$99.7\%$ by
variance), splitting into a reducible frame-tie and an irreducible Moon-excess
($\approx\!5$--$6$~nrad, a $\approx\!\SI{2}{m}$ geocentric-Moon-state envelope at the
Earth--Moon lever arm) for the JPL-versus-others pairs; this decomposition is
validated against an independent SciPy reference to a tolerance below $10^{-3}$.
Folding the two together gives the design law: with an irreducible fraction of
$0.687$ (the irreducible rotation about twice the reducible), a common
format-and-frame tag leaves the dominant $\approx\!\SI{1.8}{m}$ geocentric-Moon-state
floor---only millimetres for a surface user---and closing it requires mandating a
common (or consistency-bounded) dynamical reference.

The design law is compact enough to state in one line: \emph{for cross-provider lunar
PNT, standardize the ephemeris, not just the format.} It gives interoperability
programmes---LNIS AD-5, ESA Moonlight/LCNS, the CCSDS/IOAG cislunar standards---a
concrete, real-data-backed requirement structure and a tolerance map, on an open,
fully reproducible substrate. The theorem and the gauge lemma are proven; the
inter-ephemeris decomposition is validated on real data; the system-level budget and
tolerance map are explicitly modelled, because the providers are ephemerides and no
two lunar-PNT providers fly. This paper is the second of a programme on lunar-PNT
requirements science built on the same substrate as the companion datum-%
identifiability paper \citep{baweja2026datum}; before any absolute budget number is
trusted, a real multi-provider deployment must substitute flight data for the
representative ephemeris analog used here. The convention-free content---an
irreducible, orientation-type, Moon-specific disagreement between the real lunar
dynamical references, removable by no coordinate convention---will carry over
unchanged.
